% Luc Destombe --- licence CC BY-NC-SA 4.0
% https://creativecommons.org/licenses/by-nc-sa/4.0/deed.fr
% Fichier autonome : compiler avec pdflatex, deux fois (signets, nombre de pages).
\def\lyceeversion{1.2}
\documentclass[a4paper,10pt]{article}
\makeatletter
\newif\iflycee@niveaux
\newif\iflycee@palatino
\lycee@palatinotrue

\RequirePackage[utf8]{inputenc}
\RequirePackage[T1]{fontenc}
\RequirePackage[french]{babel}
\RequirePackage{etoolbox}

\newcommand{\ShorthandsOff}{\@ifundefined{shorthandoff}{}{\shorthandoff{;:!?}}}
\newcommand{\ShorthandsOn}{\@ifundefined{shorthandon}{}{\shorthandon{;:!?}}}
\ShorthandsOff
\AtBeginDocument{\ShorthandsOn}
\AtBeginEnvironment{tikzpicture}{\ShorthandsOff}

\RequirePackage{geometry}
\RequirePackage{fancyhdr}
\RequirePackage{lastpage}
\RequirePackage{multicol}
\RequirePackage{array}
\RequirePackage{enumitem}
\RequirePackage{xcolor}
\RequirePackage{needspace}

\RequirePackage{amsmath,amssymb}
\RequirePackage{mathpazo}
\DeclareMathSizes{10.95}{10.95}{8}{5.5}
\begingroup\catcode`\_=\active \catcode`\^=\active
\gdef_#1{\sb{\scriptscriptstyle #1}}
\gdef^#1{\sp{\scriptscriptstyle\lycee@moinsepais #1}}
\endgroup
\newcommand\lycee@moins{\mathbin{%
  \setbox\z@\hbox{$\m@th\scriptscriptstyle\lycee@moinsorig$}%
  \dimen@=.55\wd\z@ \advance\dimen@-.5pt
  \hbox to.75\wd\z@{\hss$\m@th\scriptscriptstyle\vcenter{\hbox{%
    \tikz\draw[line width=.5pt,line cap=round](0,0)--(\the\dimen@,0);}}$\hss}}}
\begingroup\lccode`\~=`\- \lowercase{\endgroup\let~}\lycee@moins
\newcommand\lycee@moinsepais{\mathcode`\-="8000 }
\AtBeginDocument{\mathchardef\lycee@moinsorig=\mathcode`\-\relax}
\AtBeginDocument{\catcode`\_=12 \mathcode`\_="8000
  \catcode`\^=12 \mathcode`\^="8000 }
\newcommand\lycee@exposants{%
  \fontdimen13\textfont2=.48\fontdimen6\textfont2
  \fontdimen14\textfont2=.48\fontdimen6\textfont2
  \fontdimen15\textfont2=.48\fontdimen6\textfont2
  \fontdimen13\scriptfont2=.48\fontdimen6\scriptfont2
  \fontdimen14\scriptfont2=.48\fontdimen6\scriptfont2
  \fontdimen15\scriptfont2=.48\fontdimen6\scriptfont2}
\every@math@size\expandafter{\the\every@math@size\lycee@exposants}
\newlength{\retraitcalculs}
\setlength{\retraitcalculs}{2em}

\RequirePackage{tikz}
\usetikzlibrary{arrows.meta,calc,babel,shapes.misc}
\RequirePackage[most]{tcolorbox}
\tcbuselibrary{skins,breakable}

\definecolor{coulDef}{HTML}{1F4E79}
\definecolor{coulProp}{HTML}{7030A0}
\definecolor{coulRem}{HTML}{C55A11}
\definecolor{coulEx}{HTML}{2E7D32}
\definecolor{coulExo}{HTML}{B03A2E}
\definecolor{coulPart}{HTML}{1F4E79}
\definecolor{coulMeth}{HTML}{1B6E4F}
\definecolor{coulCourbe}{HTML}{0B5ED7}
\definecolor{coulCourbeB}{HTML}{C00000}
\definecolor{coulCourbeC}{HTML}{2E7D32}
\definecolor{coulGrille}{HTML}{8C8C8C}
\definecolor{coulRep}{HTML}{1B6E4F}
\definecolor{coulBar}{HTML}{2E7D32}

\newcommand{\R}{\mathbb{R}}
\newcommand{\Rs}{\mathbb{R}^{*}}
\renewcommand{\le}{\leqslant}
\renewcommand{\ge}{\geqslant}
\renewcommand{\approx}{\simeq}
\newcommand{\vect}[1]{\overrightarrow{#1}}
\newcommand{\norme}[1]{\left\lVert #1 \right\rVert}
\newcommand{\intervalle}[2]{\left[#1\,;#2\right]}
\RequirePackage{cancel}
\renewcommand{\CancelColor}{\color{coulExo}}
\newcommand{\fois}[1]{\times{\color{coulExo}#1}}
\newcommand{\barre}[1]{\cancel{#1}}

\tikzset{grille/.style={coulGrille,line width=0.4pt}}
\newcommand{\quadrillage}[4]{\draw[grille,step=1] (#1,#2) grid (#3,#4);}
\tikzset{
  croix/.style={cross out,draw,line width=0.6pt,inner sep=0pt,minimum size=4pt},
  croixdroite/.style={croix,rotate=45,line width=0.8pt},
}
\newcommand{\pt}[3]{\path (#1) node[croix]{} node[#2,font=\small,inner sep=2.5pt]{$#3$};}

\newcommand{\lycee@repere}[4]{%
  \draw[grille,step=1] (#1,#2) grid (#3,#4);
  \draw[-{Stealth[length=2mm]},thick] (#1,0)--({#3+0.4},0) node[below left=-1pt]{$x$};
  \draw[-{Stealth[length=2mm]},thick] (0,#2)--(0,{#4+0.4}) node[below left=-1pt]{$y$};
  \node[below left,font=\scriptsize,inner sep=1.5pt] at (0,0) {$0$};}
\newcommand{\lycee@gradx}[1]{\foreach \k/\l in {#1}{%
  \draw (\k,0.07)--(\k,-0.07) node[below,font=\scriptsize,inner sep=1.5pt]{$\l$};}}
\newcommand{\lycee@grady}[1]{\foreach \k/\l in {#1}{%
  \draw (0.07,\k)--(-0.07,\k) node[left,font=\scriptsize,inner sep=1.5pt]{$\l$};}}
\newcommand{\lycee@intff}[2]{\left[#1\,;#2\right]}
\newcommand{\lycee@intfo}[2]{\left[#1\,;#2\right[}
\newcommand{\lycee@intof}[2]{\left]#1\,;#2\right]}
\newcommand{\lycee@intoo}[2]{\left]#1\,;#2\right[}
\newcommand{\lycee@ens}[1]{\left\{#1\right\}}
\AtBeginDocument{%
  \providecommand{\repere}{\lycee@repere}%
  \providecommand{\gradx}{\lycee@gradx}%
  \providecommand{\grady}{\lycee@grady}%
  \providecommand{\Cf}{\mathcal{C}\sb{\scriptscriptstyle f}}%
  \providecommand{\Cg}{\mathcal{C}\sb{\scriptscriptstyle g}}%
  \providecommand{\intff}{\lycee@intff}%
  \providecommand{\intfo}{\lycee@intfo}%
  \providecommand{\intof}{\lycee@intof}%
  \providecommand{\intoo}{\lycee@intoo}%
  \providecommand{\ens}{\lycee@ens}}

\newlist{questions}{enumerate}{3}
\setlist[questions,1]{label=\arabic*\textdegree),leftmargin=*,itemsep=2pt,topsep=3pt}
\setlist[questions,2]{label=\alph*),leftmargin=*,itemsep=1pt,topsep=2pt}
\setlist[questions,3]{label=\roman*),leftmargin=*,itemsep=1pt,topsep=1pt}
\setlist[itemize]{label=\textbullet,leftmargin=*,itemsep=2pt,topsep=3pt}

\newcommand{\besoinplace}[1]{\par\penalty\@M\begingroup
  \dimen@=#1\relax\dimen@ii=\pagegoal\advance\dimen@ii-\pagetotal
  \ifdim\dimen@>\dimen@ii\ifdim\dimen@ii>\z@\vfil\fi\break\fi\endgroup}

\newcommand{\lycee@pied}{\thepage/\pageref{LastPage}}
\newcommand{\entete}[2]{%
  \pagestyle{fancy}\fancyhf{}%
  \fancyhead[L]{\footnotesize\color{coulPart}\textbf{#1}}%
  \fancyhead[R]{\footnotesize\color{coulPart}\textbf{#2}}%
  \fancyfoot[C]{\footnotesize\lycee@pied}%
  \renewcommand{\headrulewidth}{0.4pt}%
  \renewcommand{\headrule}{\hbox to\headwidth{%
    \color{coulPart!40}\leaders\hrule height \headrulewidth\hfill}}}

\RequirePackage{ccicons}
\newcommand{\lycee@licence}{{\scriptsize\color{gray}%
  \copyright\ Luc Destombe --- \ccbyncsa\ CC BY-NC-SA 4.0}}
\newif\iflycee@licence \lycee@licencetrue
\newcommand{\sanslicence}{\lycee@licencefalse}
\AtBeginDocument{\iflycee@licence\AddToHookNext{shipout/foreground}{%
  \put(0,0){\rlap{\kern\dimexpr1in+\hoffset+\oddsidemargin+\textwidth\relax
    \llap{\raisebox{-\dimexpr1in+\voffset+\topmargin+\headheight+\headsep
      +\textheight+\footskip\relax}{\lycee@licence}}}}}\fi}

\newcommand{\lycee@site}{faire-des-maths.fr}
\newcommand{\lycee@version}{{\scriptsize\color{gray}\lycee@site
  \ifdefined\lyceeversion\ --- Version~\lyceeversion\fi}}
\AtBeginDocument{\iflycee@licence\AddToHookNext{shipout/foreground}{%
  \put(0,0){\rlap{\kern\dimexpr1in+\hoffset+\oddsidemargin\relax
    \raisebox{-\dimexpr1in+\voffset+\topmargin+\headheight+\headsep
      +\textheight+\footskip\relax}{\lycee@version}}}}\fi}

\newcommand{\formulesserrees}{%
  \setlength{\abovedisplayskip}{3pt}\setlength{\belowdisplayskip}{3pt}%
  \setlength{\abovedisplayshortskip}{2pt}\setlength{\belowdisplayshortskip}{3pt}}

\setlength{\parindent}{0pt}

\RequirePackage[implicit=false,hidelinks,pdfencoding=auto,psdextra,bookmarksopen,
  bookmarksopenlevel=1,pdfcreator={},pdfproducer={}]{hyperref}
\RequirePackage{bookmark}
\hypersetup{pdfauthor={Luc Destombe},pdfkeywords={CC BY-NC-SA 4.0}}
\newcounter{lycee@signet}
\newcommand{\signet}[3][]{\stepcounter{lycee@signet}%
  \edef\lycee@dest{\if\relax\detokenize{#1}\relax
    signet.\arabic{lycee@signet}\else#1\fi}%
  \hypertarget{\lycee@dest}{}\bookmark[level=#2,dest=\lycee@dest]{#3}}
\pdfstringdefDisableCommands{%
  \def\R{R}\def\vect#1{#1}\def\Cf{Cf}\def\Cg{Cg}\def\fois#1{×#1}%
  \def\barre#1{#1}\def\intervalle#1#2{[#1;#2]}\def\intff#1#2{[#1;#2]}%
  \def\textsuperscript#1{#1}\def\dfrac#1#2{#1/#2}\def\frac#1#2{#1/#2}%
  \def\vec#1{#1⃗}}%
\RequirePackage[official]{eurosym}

\geometry{top=0.8cm,bottom=1.3cm,left=1cm,right=1cm,footskip=0.6cm}

\pagestyle{fancy}
\fancyhf{}
\renewcommand{\headrulewidth}{0pt}
\fancyfoot[C]{\footnotesize Page \thepage{} / \pageref{LastPage}}

\newcommand{\titrefiche}[2]{%
  \begin{center}
    \Large\textbf{#1}\\[2pt]
    \large\textbf{#2}\\[2pt]
  \end{center}
  \vspace{0.10cm}}

\setlength{\columnseprule}{0.4pt}
\setlength{\columnsep}{15pt}
\renewcommand{\columnseprulecolor}{\color{black!45}}
\setlength{\multicolsep}{2pt}

\newcounter{partiefiche}
\renewcommand{\thepartiefiche}{\Roman{partiefiche}}
\newif\ifapresbandeau
\newcommand{\lycee@bandeau}[1]{%
  \par\penalty-600
  \vspace{0.08cm}%
  \noindent\colorbox{coulPart}{%
    \parbox{\dimexpr\linewidth-2\fboxsep\relax}{%
      \color{white}\bfseries\raggedright\strut#1\strut}}%
  \nopagebreak\par\nopagebreak
  \vspace{0.06cm}\nopagebreak
  \apresbandeautrue}
\newcommand{\partiefiche}[1]{%
  \refstepcounter{partiefiche}%
  \lycee@bandeau{\signet{1}{\thepartiefiche{} --- #1}\thepartiefiche{} --- #1}}
\newcommand{\partiefichesansnum}[1]{\lycee@bandeau{\signet{1}{#1}#1}}

\newcounter{exo}
\newlength{\espaceexo}\setlength{\espaceexo}{7pt}
\newlength{\espacetitre}\setlength{\espacetitre}{2pt}
\newcommand{\exo}[1][\relax]{%
  \par
  \ifapresbandeau\apresbandeaufalse\else\penalty-150\addvspace{\espaceexo}\fi
  \refstepcounter{exo}%
  \noindent\signet[exercice-\arabic{exo}]{2}{Exercice \arabic{exo}}%
  \textbf{\color{coulExo}\underline{Exercice \theexo}}%
  \ifx\relax#1\relax\else\textbf{ : #1}\fi
  \par\nopagebreak\vspace{\espacetitre}\nopagebreak
  \ignorespaces}

\setlist[enumerate]{nosep,leftmargin=*,itemsep=2pt,topsep=2pt}
\setlist[itemize]{nosep,leftmargin=*,topsep=2pt}
\newcommand{\choix}[3]{\par\hspace*{1em}%
  \makebox[0.3\linewidth][l]{a) #1}\makebox[0.3\linewidth][l]{b) #2}c) #3}
\newcommand{\deux}[2]{\par\hspace*{0.6em}%
  \makebox[0.48\linewidth][l]{#1}#2}
\newcommand{\trois}[3]{\par\hspace*{0.6em}%
  \makebox[0.32\linewidth][l]{#1}\makebox[0.32\linewidth][l]{#2}#3}

\renewenvironment{center}{\par\addvspace{3pt}\centering}{\par\addvspace{3pt}}
\formulesserrees
\AtBeginDocument{\formulesserrees}
\clubpenalty=10000
\widowpenalty=10000
\displaywidowpenalty=10000
\brokenpenalty=10000
\setlength{\parskip}{0pt}

\RequirePackage{ragged2e}
\setlength{\RaggedRightRightskip}{0pt plus 1fil}
\AtBeginDocument{\RaggedRight\binoppenalty=10000 \relpenalty=10000 }
\g@addto@macro\@parboxrestore{\RaggedRight}
\makeatother
\geometry{bottom=1cm}

\setlength{\espaceexo}{4pt}
\setlist[enumerate]{nosep,leftmargin=*,itemsep=1pt,topsep=1pt}

\begin{document}

\titrefiche{Seconde --- Pourcentages}{Fiche d'exercices du chapitre}

\begin{multicols}{2}\raggedcolumns

\partiefiche{Appliquer et calculer un pourcentage}

\exo[Prendre un pourcentage]
Calculer, à la main ou à la calculatrice :
\begin{enumerate}[label=\alph*)]
  \item $10\,\%$ de $340$ ;\; $25\,\%$ de $48$ ;\; $50\,\%$ de $3\,600$ ;
  \item $15\,\%$ de $260$ ;\; $8\,\%$ de $75$ ;\; $3{,}5\,\%$ de $400$ ;
  \item $200\,\%$ de $17$ ;\; $120\,\%$ de $45$ ;\; $0{,}5\,\%$ de $1\,800$.
\end{enumerate}

\exo[Appliquer un pourcentage]
\begin{enumerate}[label=\alph*)]
  \item Un lycée compte $1\,150$ élèves, dont $42\,\%$ sont demi-pensionnaires.
        Combien d'élèves sont demi-pensionnaires ?
  \item Un téléphone coûte $420$~€ ; il a une réduction de $15\,\%$. Calculer
        le montant de la réduction, puis le prix payé.
  \item Le téléchargement d'un fichier de $360$~Mo en est à $65\,\%$. Combien
        de Mo reste-t-il à télécharger ?
\end{enumerate}

\exo[Calculer un pourcentage]
\begin{enumerate}[label=\alph*)]
  \item À une question d'un QCM, $18$ élèves ont répondu juste et $7$ ont
        répondu faux. Quel pourcentage des élèves a répondu juste ?
  \item Une entreprise de $320$ salariés embauche $40$ personnes. Quel
        pourcentage de l'effectif de départ représentent les embauches ?
  \item Sur un salaire de $1\,600$~€, Malo dépense $288$~€ en courses. Quel
        pourcentage de son salaire représentent les courses ?
\end{enumerate}

\exo[Comparer des pourcentages]
Dans la classe A, $12$ élèves sur $30$ suivent l'option théâtre ; dans la classe
B, $13$ sur $35$ ; dans la classe C, $11$ sur $28$.
\begin{enumerate}[label=\alph*)]
  \item Pour chaque classe A, B et C, calculer le pourcentage d'élèves qui
        suivent l'option (arrondir au dixième).
  \item Ranger les classes de la plus forte part d'élèves de théâtre à la plus
        faible. La classe qui a le plus d'élèves de théâtre est-elle la
        première ?
\end{enumerate}

\exo[Lire un tableau]
Les $1\,200$ élèves d'un lycée ont indiqué comment ils viennent au lycée.
\begin{center}
\begin{tabular}{|l|c|c|c|c|c|}
\hline
Mode & bus & vélo & à pied & voiture & total\\
\hline
Élèves & $504$ & $156$ & $300$ & $240$ & $1\,200$\\
\hline
\end{tabular}
\end{center}
\begin{enumerate}[label=\alph*)]
  \item Pour chaque mode de transport, calculer le pourcentage d'élèves qui
        l'utilisent.
  \item Écrire une phrase qui explique ce que représente le quotient
        $\dfrac{300}{1\,200}$ dans ce contexte.
\end{enumerate}

\exo[La TVA]
La TVA est une taxe de $20\,\%$ ajoutée au prix hors taxe (HT) ; on obtient le
prix toutes taxes comprises (TTC).
\begin{enumerate}[label=\alph*)]
  \item Un bureau coûte $240$~€ HT. Calculer le montant de la TVA, puis le prix
        TTC.
  \item Même question pour une chaise à $85$~€ HT.
\end{enumerate}

\partiefiche{Proportions}

\exo[Effectif, proportion, effectif total]
Un club de judo compte $320$ licenciés, dont $35\,\%$ sont mineurs. Le club
compte $144$ licenciées, dont $72$ jeunes filles (mineures).
\begin{enumerate}[label=\alph*)]
  \item Combien de licenciés sont mineurs ?
  \item Quelle est la proportion de licenciées dans tout le club ?
  \item Quelle est la proportion de jeunes filles parmi les licenciées ? La
        comparer à la proportion de mineurs dans tout le club.
  \item Un autre club compte $54$ licenciés mineurs, soit $30\,\%$ de ses
        licenciés. Combien ce club a-t-il de licenciés ?
\end{enumerate}

\exo[Population de référence]
Un constructeur vend $2{,}4$ millions de voitures en un an, dont $360\,000$
voitures électriques.
\begin{enumerate}[label=\alph*)]
  \item Quelle est la population de référence ? Quelle est la sous-population ?
        Répondre par deux phrases.
  \item Calculer la proportion de voitures électriques, en pourcentage.
\end{enumerate}

\exo[Retrouver l'effectif total]
\begin{enumerate}[label=\alph*)]
  \item Dans une population d'insectes, $300$ insectes portent un gène rare ;
        ils représentent $1{,}5\,\%$ de la population. Combien la population
        compte-t-elle d'insectes ?
  \item Un réseau social annonce $4{,}8$ milliards d'utilisateurs, soit $60\,\%$
        de la population mondiale. Quelle population mondiale cela suppose-t-il ?
  \item Une classe compte $21$ filles, soit $60\,\%$ des élèves. Combien la
        classe compte-t-elle d'élèves ?
\end{enumerate}

\exo[Un budget]
Le loyer de Yanis, $640$~€, représente $40\,\%$ de son salaire.
\begin{enumerate}[label=\alph*)]
  \item Calculer le salaire de Yanis.
  \item Ses charges représentent $8\,\%$ de son salaire. Calculer leur montant.
  \item Quel pourcentage de son salaire représentent le loyer et les charges
        ensemble ?
\end{enumerate}

\exo[Vérifier une information]
Un journal affirme : « $1{,}2$ million de personnes ont eu la grippe cet hiver,
soit $3\,\%$ de la population. »
\begin{enumerate}[label=\alph*)]
  \item Si l'affirmation était juste, quelle serait la population du pays ?
  \item Le pays compte en réalité $68$ millions d'habitants. Quel pourcentage de
        la population a eu la grippe (arrondir au dixième) ?
\end{enumerate}

\partiefiche{Proportions échelonnées}

\exo[Pourcentage d'un pourcentage]
Calculer :\; $20\,\%$ de $40\,\%$ ;\; $50\,\%$ de $30\,\%$ ;\; $10\,\%$ de $5\,\%$.

\exo[Association sportive]
Dans un lycée de $1\,200$ élèves, $45\,\%$ des élèves sont inscrits à
l'association sportive, et $20\,\%$ des inscrits font du badminton.
\begin{enumerate}[label=\alph*)]
  \item Quelle proportion des élèves du lycée fait du badminton à
        l'association sportive ?
  \item Combien d'élèves cela représente-t-il ?
  \item Vérifier ce résultat en calculant d'abord le nombre d'inscrits.
\end{enumerate}

\exo[Populations emboîtées]
\begin{enumerate}[label=\alph*)]
  \item Une compote contient $80\,\%$ de fruits, dont $35\,\%$ de pommes. Quel
        est le pourcentage de pommes dans la compote ?
  \item Dans un sachet de $300$ bonbons, $25\,\%$ sont à la fraise, et $40\,\%$
        de ceux-ci sont acidulés. Combien le sachet contient-il de bonbons
        acidulés à la fraise ?
  \item Une ville compte $48\,000$ habitants ; $75\,\%$ sont majeurs, $60\,\%$
        des majeurs ont le permis de conduire, et $25\,\%$ de ces derniers
        roulent en voiture électrique. Quelle proportion des habitants roule en
        voiture électrique ? Combien d'habitants cela représente-t-il ?
\end{enumerate}

\exo[Tableau croisé]
Un collège compte $600$ élèves ; $55\,\%$ sont des filles, et $90\,\%$ des
filles sont droitières. Parmi les garçons, $10\,\%$ sont gauchers.
\begin{enumerate}[label=\alph*)]
  \item Recopier et compléter le tableau d'effectifs.
\begin{center}
\begin{tabular}{|l|c|c|c|}
\hline
 & Gauchers & Droitiers & Total\\
\hline
Filles & & & \\
\hline
Garçons & & & \\
\hline
Total & & & $600$\\
\hline
\end{tabular}
\end{center}
  \item Quelle est la proportion de filles gauchères dans le collège ? La
        retrouver par une proportion échelonnée.
  \item Quelle est la proportion de gauchers dans le collège ?
\end{enumerate}

\partiefiche{Taux d'évolution}

\exo[Taux d'évolution]
Pour chaque grandeur, calculer la variation absolue, puis le taux d'évolution
en pourcentage.
\begin{enumerate}[label=\alph*)]
  \item Une commune passe de $3\,200$ à $3\,040$ habitants.
  \item Une usine produit $1\,500$ vélos une année, puis $2\,100$ l'année
        suivante.
  \item Un salaire passe de $1\,650$~€ à $1\,782$~€.
  \item Un plein de carburant passe de $18{,}40$~€ à $19{,}10$~€ (arrondir au
        dixième de pourcentage).
\end{enumerate}

\exo[Deux façons de dire]
Un club passe de $400$ à $460$ adhérents.
\begin{enumerate}[label=\alph*)]
  \item Calculer la variation absolue et le taux d'évolution.
  \item Écrire une phrase qui donne la variation absolue, puis une phrase qui
        donne le taux d'évolution.
\end{enumerate}

\exo[Deux magasins]
Un même téléviseur est vendu dans deux magasins. Dans le magasin A, son prix
passe de $480$~€ à $408$~€. Dans le magasin B, il coûtait $500$~€ et baisse de
$14\,\%$.
\begin{enumerate}[label=\alph*)]
  \item Dans quel magasin la baisse est-elle la plus forte, en pourcentage ?
  \item Dans quel magasin le téléviseur est-il le moins cher ?
\end{enumerate}

\exo[Une entreprise qui embauche]
Une entreprise compte $250$ salariés. Elle embauche $30$ personnes une année,
puis $14$ l'année suivante.
\begin{enumerate}[label=\alph*)]
  \item Combien l'entreprise compte-t-elle de salariés après la première
        année ?
  \item Quel pourcentage de l'effectif de départ représentent les $30$
        embauches ?
  \item Combien l'entreprise compte-t-elle de salariés après la seconde année ?
  \item Calculer le taux d'évolution du nombre de salariés sur les deux années.
\end{enumerate}

\partiefiche{Coefficient multiplicateur}

\exo[Taux et coefficient multiplicateur]
\begin{enumerate}[label=\alph*)]
  \item Pour chaque évolution, donner le coefficient multiplicateur $C_m$ :
        hausse de $15\,\%$ ;\; baisse de $4\,\%$ ;\; hausse de $120\,\%$ ;\;
        baisse de $60\,\%$ ;\; baisse de $12{,}3\,\%$.
  \item Pour chaque coefficient multiplicateur, donner le taux d'évolution :
        $C_{m1}=1{,}03$ ;\; $C_{m2}=0{,}72$ ;\; $C_{m3}=2$ ;\; $C_{m4}=0{,}995$.
\end{enumerate}

\exo[En phrases]
Pour chaque coefficient multiplicateur, écrire une phrase qui décrit
l'évolution, sur le modèle : « Multiplier par $1{,}25$, c'est augmenter de
$25\,\%$. »
\[
  1{,}08 \qquad 0{,}92 \qquad 1{,}5 \qquad 0{,}3
\]

\exo[Du coefficient au taux]
Une valeur passe de $V_i$ à $V_f$. On a $V_f=V_i\times C_m$, donc
$C_m=\dfrac{V_f}{V_i}$, et le taux d'évolution est $t=C_m-1$.
\begin{enumerate}[label=\alph*)]
  \item Un prix passe de $120$~€ à $144$~€. Calculer $C_m$, puis $t$.
  \item Un stock passe de $250$~kg à $200$~kg. Calculer $C_m$, puis $t$.
  \item Un prix est multiplié par $1{,}5$. Calculer $t$.
  \item Un prix est multiplié par $0{,}65$. Calculer $t$.
\end{enumerate}

\exo[Vrai ou faux]
Pour chaque affirmation, dire si elle est vraie ou fausse, et justifier.
\begin{enumerate}[label=\alph*)]
  \item Un prix passe de $120$~€ à $150$~€ : il augmente de $30\,\%$.
  \item Un prix passe de $120$~€ à $150$~€ : il est multiplié par $1{,}25$.
  \item Un stock passe de $250$~kg à $200$~kg : il baisse de $25\,\%$.
  \item Un village de $300$ habitants en perd $60$ : sa population baisse de
        $20\,\%$.
\end{enumerate}

\exo[Valeur finale]
Pour chaque situation, calculer la nouvelle valeur avec le coefficient
multiplicateur.
\begin{enumerate}[label=\alph*)]
  \item Une ville de $24\,000$ habitants gagne $5\,\%$ d'habitants.
  \item Des chaussures à $75$~€ ont une réduction de $20\,\%$.
  \item Dans un département, $3\,500$ accidents de la route ont eu lieu une
        année ; leur nombre baisse de $12\,\%$ l'année suivante.
\end{enumerate}

\exo[Valeur initiale]
\begin{enumerate}[label=\alph*)]
  \item Après une hausse de $8\,\%$, un abonnement coûte $27$~€. Quel était son
        prix avant la hausse ?
  \item En septembre, la température moyenne d'une ville est de $19$~°C, en
        baisse de $5\,\%$ par rapport à août. Quelle était la température
        moyenne en août ?
  \item Un lit coûte $150$~€ TTC, avec une TVA de $20\,\%$. Quel est son prix
        HT ?
  \item Après une hausse de $18\,\%$, le prix d'un vélo a augmenté de $27$~€.
        Quel était son prix avant la hausse ?
\end{enumerate}

\exo[Les soldes]
Pendant les soldes, un magasin fait une réduction de $20\,\%$ sur tous ses
articles.
\begin{enumerate}[label=\alph*)]
  \item Léo a dépensé $225$~€ juste avant les soldes. Combien aurait-il payé
        pendant les soldes ?
  \item Inès a dépensé $240$~€ pendant les soldes. Combien aurait-elle payé
        avant les soldes ?
\end{enumerate}

\partiefiche{Évolutions successives}

\exo[Taux global]
Pour chaque cas, calculer le coefficient multiplicateur global, puis le taux
d'évolution global.
\begin{enumerate}[label=\alph*)]
  \item Hausse de $25\,\%$, puis hausse de $10\,\%$.
  \item Hausse de $30\,\%$, puis baisse de $30\,\%$.
  \item Baisse de $40\,\%$, puis baisse de $10\,\%$.
  \item Hausse de $2\,\%$, puis baisse de $1\,\%$.
\end{enumerate}

\exo[Un ticket de transport]
Le prix d'un abonnement de transport augmente de $30\,\%$ ; jugée trop forte,
la hausse est suivie d'une baisse de $10\,\%$.
\begin{enumerate}[label=\alph*)]
  \item Quel est le taux d'évolution global ?
  \item L'abonnement coûtait $25$~€. Combien coûte-t-il après les deux
        évolutions ?
\end{enumerate}

\exo[Vrai ou faux]
Pour chaque affirmation, dire si elle est vraie ou fausse, et justifier.
\begin{enumerate}[label=\alph*)]
  \item Une hausse de $12\,\%$ suivie d'une baisse de $12\,\%$ ne change pas le
        prix.
  \item Une réduction de $7\,\%$ puis de $20\,\%$ donne le même prix qu'une
        réduction de $20\,\%$ puis de $7\,\%$.
  \item Une hausse de $6\,\%$ puis de $8\,\%$ fait une hausse de $14\,\%$.
  \item « Une hausse de $10\,\%$ en mars, puis de $15\,\%$ en septembre : le prix
        a augmenté d'un quart sur l'année. »
\end{enumerate}

\exo[Le fleuriste]
Un fleuriste achète des orchidées $8$~€ le pot. Il augmente ce prix de
$40\,\%$ pour les vendre, puis, pour les fêtes, fait une réduction de $30\,\%$
sur son prix de vente. Vend-il ses orchidées à perte ?

\exo[La valeur d'une voiture]
Une voiture neuve coûte $18\,000$~€. Elle perd $20\,\%$ de sa valeur la
première année, puis $15\,\%$ par an les années suivantes.
\begin{enumerate}[label=\alph*)]
  \item Calculer sa valeur au bout d'un an, puis de deux ans.
  \item Calculer sa valeur au bout de trois ans.
  \item Quel est le taux d'évolution global sur les trois ans (arrondir au
        dixième de pourcentage) ?
\end{enumerate}

\partiefiche{Évolution réciproque}

\exo[Revenir au départ]
Pour chaque évolution, quelle évolution ramène la valeur à son point de
départ ? Arrondir au dixième de pourcentage si besoin.
\begin{enumerate}[label=\alph*)]
  \item Baisse de $20\,\%$.
  \item Hausse de $100\,\%$.
  \item Baisse de $45\,\%$.
  \item Hausse de $30\,\%$.
\end{enumerate}

\exo[Un village]
La population d'un village passe de $2\,500$ à $3\,000$ habitants en un an.
\begin{enumerate}[label=\alph*)]
  \item Calculer le taux d'évolution de la population.
  \item Quelle baisse, en pourcentage, ramènerait le village à $2\,500$
        habitants (arrondir au dixième) ?
\end{enumerate}

\exo[Aller et retour]
Pour chaque prix, calculer le taux d'évolution, le coefficient multiplicateur
$C_m$, puis le taux de l'évolution qui ramène au prix de départ (arrondir au
dixième de pourcentage).
\begin{enumerate}[label=\alph*)]
  \item Un prix passe de $20$~€ à $26$~€.
  \item Un prix passe de $52$~€ à $39$~€.
\end{enumerate}

\exo[Plus de chocolat pour le même prix]
Pour le même prix, une tablette de chocolat contient $25\,\%$ de chocolat en
plus. De quel pourcentage le prix du kilo de chocolat a-t-il baissé ? On pourra
prendre un exemple : $10$~€ pour $1$~kg.

\exo[Retrouver un prix]
Un casque a subi une hausse de $10\,\%$, puis une baisse de $20\,\%$ ; il coûte
maintenant $352$~€.
\begin{enumerate}[label=\alph*)]
  \item Calculer le coefficient multiplicateur global et le taux d'évolution
        global.
  \item Quel était le prix du casque avant les deux évolutions ?
  \item Quelle évolution ramènerait le prix à sa valeur initiale (arrondir au
        dixième de pourcentage) ?
\end{enumerate}

\end{multicols}

\end{document}
